\documentclass[../../../main]{subfiles}

\ifSubfilesClassLoaded{%
    \zexternaldocument*{../../1.Category-Theory/exported/aux/main}%
    \zexternaldocument*{../../2.Homological-Algebra/exported/aux/main}%
    \zexternaldocument*{../../3.Sheaves/1.Sheaves/exported/aux/main}%
    \zexternaldocument*{./1.Sheaves-of-modules}%
    \zexternaldocument*{./2.Ringed-spaces}%
}{}

\begin{document}
\section{Vector bundles}\label{sec:vector-bundles}

\begin{definition}[Topological vector bundle]\label{def:topological-vector-bundle}
  Let \(X\) be a topological space.
  Let \(k \in \{\mathbb{R}, \mathbb{C}\}\) be either the field of real or complex numbers.
 
  A \textbf{topological \(k\)-vector bundle} over \(X\) (denoted by \(E \xrightarrow{\pi} X\)) consists of
  \begin{itemize}
    \item A topological space \(E\).
    \item A continuous map \(\pi: E \to X\).
    \item For every point \(x \in X\) the structure of a finite-dimensional \(k\)-vector space over \(E_{x} \coloneq \pi^{-1}(\{x\})\).
  \end{itemize}
  That admits a \textit{local trivialization}, i.e: there exists an open cover \(\{U_{i}\}_{i \in \mathcal{I}}\) of \(X\) and for each \(i \in \mathcal{I}\) an \(n_{i} \in \mathbb{N}\) and a homeomorphism
  \begin{equation*}
    \phi_{i}: U_{i} \times k^{n_{i}} \isoarrow \pi^{-1}(U_{i})
  \end{equation*}
  satisfying
  \begin{enumerate}
    \item \(\pi \circ \phi_{i} = \operatorname{pr}_{1}\), meaning that \(\phi_{i}(\{x\} \times k^{n_{i}}) \subseteq \pi^{-1}(\{x\})\).
    \item For all \(x \in X\) the map \(\phi_{i}(x): k^{n_{i}} \to E_{x}\) is linear.
  \end{enumerate}
  If all the numbers \(n_{i}\) are equal to some \(n \in \mathbb{N}\) we say that the vector bundle has rank \(n\).
  
  A \textbf{morphism} of topological \(k\)-vector bundles from \(E_{1} \xrightarrow{\pi_{1}} X\) to \(E_{2} \xrightarrow{\pi_{2}} X\) is a continuous map \(f: E_{1} \to E_{2}\) such that \(\pi_{2} \circ f = \pi_{1}\) and for every \(x \in X\) the map \(f|_{x}: (E_{1})_{x} \to (E_{2})_{x}\) is linear.

  We denote by \(\Vect_{k}(X)\) the category of topological \(k\)-vector bundles over \(X\).
\end{definition}

\begin{notation}
  A \hyperref[def:topological-vector-bundle]{topological \(k\)-vector bundle} of rank \(n\) is:
  \begin{itemize}
    \item A \textbf{real topological vector bundle} if \(k = \mathbb{R}\).
    \item A \textbf{complex topological vector bundle} if \(k = \mathbb{C}\).
    \item A \textbf{line bundle} if \(n = 1\).
  \end{itemize}
\end{notation}


\begin{example}[Trivial topological vector bundle]\label{ex:trivial-topological-vector-bundle}
  Let \(X\) be a topological space and \(n \in \mathbb{N}\).
  The product topological space \(X \times k^{n} \xrightarrow{\operatorname{pr}_{1}} X\) is a \hyperref[def:topological-vector-bundle]{topological \(k\)-vector bundle} of rank \(n\) called the \textbf{trivial} topological vector bundle over \(X\).
\end{example}

\begin{remark}
  Let \(E \xrightarrow{\pi} X\) be a \hyperref[def:topological-vector-bundle]{topological \(k\)-vector bundle} of rank \(n\) with local trivialization
  \begin{equation*}
    \{\phi_{i}: U_{i} \times k^{n} \isoarrow \pi^{-1}(U_{i})\}_{i \in \mathcal{I}}
  \end{equation*}
  For each \(i, j \in \mathcal{I}\) the composite function
  \begin{equation*}
    \phi_{i}^{-1} \circ \phi_{j}: (U_{i} \cap U_{j}) \times k^{n} \to (U_{i} \cap U_{j}) \times k^{n}
  \end{equation*}
  is well-defined and satisfies
  \begin{equation*}
    \left(  \phi_{i}^{-1} \circ \phi_{j} \right)(x, y) = (x, g_{i,j}(x)y)
  \end{equation*}
  for some \(\operatorname{GL}_{n}(k)\)-valued function
  \begin{equation*}
    g_{i,j}: U_{i} \cap U_{j} \to \operatorname{GL}_{n}(k)
  \end{equation*}
  We call these functions \(g_{i, j}\) the \textbf{transition functions} of the bundle.

  Note that this functions satisfy \(g_{i,i}(x) = \operatorname{Id}_{n}\) for all \(x \in U_{i}\) and, for all open subsets \(U_{i}, U_{j}, U_{k}\) such that the bundle \hyperref[def:trivial-vector-bundle]{trivializes} and \(U_{i} \cap U_{j} \cap U_{k} \neq \varnothing\)
  \begin{equation*}
    \forall x \in U_{i} \cap U_{j} \cap U_{k} \quad g_{i,j}(x)g_{j,k}(x)g_{k,i}(x) = \operatorname{Id}_{n}
  \end{equation*}
  This is called a \textit{\v{C}ech cocycle}.
\end{remark}

\begin{lemma}
  Let \((E, X, \pi, k^{n})\) be a fiber bundle with a \(\operatorname{GL}_{n}(k)\) cocycle acting in the standard way on the fiber \(k^{n}\).
  Then there exists a \hyperref[def:topological-vector-bundle]{topological \(k\)-vector bundle} with transition functions determined by the cocycle.
\end{lemma}


\begin{example}
  \todo[inline]{Examples of vector bundles: Cylinder, Möbius, ...}
\end{example}

\begin{definition}[Section of a topological vector bundle]\label{def:section-topological-vector-bundle}
  Let \(E \xrightarrow{\pi} X\) be a \hyperref[def:topological-vector-bundle]{topological \(k\)-vector bundle}.

  For any open set \(U \subseteq X\) a \textbf{section} of \(E \xrightarrow{\pi} X\) over \(U\) is a continuous map \(\sigma: U \to E\) such that \(\pi \circ \sigma = \id_{U}\).

  The set of all sections of \(E \xrightarrow{\pi} X\) over \(U\) is denoted by \(\Gamma(U, E \xrightarrow{\pi} X)\).
\end{definition}
\begin{remark}
  The set \(\Gamma(U, E \xrightarrow{\pi} X)\) admits the structure of a \(k\)-vector space.
\end{remark}

\begin{definition}\label{def:sheaf-of-sections-topological-vector-bundle}
  Let \(E \xrightarrow{\pi} X\) be a \hyperref[def:topological-vector-bundle]{topological \(k\)-vector bundle}.
  The \textbf{\hyperref[def:sheaf]{sheaf} of sections} of \(E \xrightarrow{\pi} X\) is
  \begin{align*}
    \mathcal{E}:\,& \Open^{\op}(X) \to \Vect_{k} \\
        &U \mapsto \Gamma(U, E \xrightarrow{\pi} X)
  \end{align*}
\end{definition}

\begin{proposition}\label{prop:sheaf-of-sections-topological-vector-bundle-is-C-x-module}
  Let \(E \xrightarrow{\pi} X\) be a \hyperref[def:topological-vector-bundle]{topological \(k\)-vector bundle}.
  Let \(\mathcal{C}_{X}(k)\) be the \hyperref[def:sheaf]{sheaf} of continuous \(k\)-valued functions over \(X\).

  Then the \hyperref[def:sheaf-of-sections-topological-vector-bundle]{sheaf of sections} \(\mathcal{E}\) is a \hyperref[def:sheaf-of-modules]{sheaf of \(\mathcal{C}_{X}(k)\)-modules}.
\end{proposition}
\begin{proof}
  Let \(U \subseteq X\) be an open subset.
  Let \(f \in \sheaf{O}_{X}(U)\) and \(\sigma \in \mathcal{E}(U)\).
  
  We define the action \(f \cdot \sigma\) pointwise: for every \(x \in U\)
  \begin{equation*}
    (f \cdot \sigma)(x) \coloneq f(x) \cdot \sigma(x)
  \end{equation*}
  where the operation in the right is scalar multiplication in \(E_{x}\).

  Let's verify that \(f \cdot \sigma\) is a valid \hyperref[def:section-topological-vector-bundle]{section} of \(E \xrightarrow{\pi} X\).
  \begin{enumerate}
      \item \textbf{Linearity:} Since \((f \cdot \sigma)(x)\) is a scalar multiple of a vector in \(E_{x}\), it remains in the fiber \(\pi^{-1}(\{x\})\). Thus \(\pi \circ (f \cdot \sigma) = \id_{U}\).
      \item \textbf{Continuity:} Since the bundle admits local trivializations, scalar multiplication \(k \times E \to E\) is a continuous map. Because \(f\) and \(\sigma\) are both continuous, their pointwise product \(f \cdot \sigma\) is continuous.
  \end{enumerate}

  The module axioms follow directly from the vector space structure on \(E_{p}\).
\end{proof}

\begin{definition}[Locally free sheaf]\label{def:locally-free-sheaf}
  A \hyperref[def:sheaf-of-modules]{sheaf of \(\sheaf{O}_{X}\)-modules} \(\sheaf{F}\) is \textbf{locally free} of rank \(n\) if there exist an open cover \(\{U_{i}\}_{i \in \mathcal{I}}\) of \(X\) such that for every \(i \in \mathcal{I}\) there is an isomorphism of \(\sheaf{O}_{X}|_{U_{i}}\)-modules
  \begin{equation*}
    \sheaf{F}|_{U_{i}} \cong \bigoplus_{j=1}^{n} \sheaf{O}_{X}|_{U_{i}} \eqcolon (\sheaf{O}_{X}|_{U_{i}})^{\oplus n}
  \end{equation*}

  We denote by \(\category{LocFree}_{n}(\sheaf{O}_{X})\) the \hyperref[def:subcategory]{full subcategory} of \(\Mod(\sheaf{O}_{X})\) consisting of locally free \hyperref[def:sheaf-of-modules]{sheaves of \(\sheaf{O}_{X}\)-modules}.
\end{definition}

\begin{theorem}\label{thm:topological-vector-bundle-is-free-sheaf}
  Let \(k \in \{\mathbb{R}, \mathbb{C}\}\).
  There is an \hyperref[def:equivalent-categories]{equivalence of categories} between the category of \hyperref[def:topological-vector-bundle]{topological \(k\)-vector bundles} of rank \(n\) over \(X\) and the category of \hyperref[def:locally-free-sheaf]{locally free sheaves of \(\mathcal{C}_{X}(k)\)-modules} of rank \(n\).

  Under this correspondence
  \begin{itemize}
    \item The \hyperref[ex:trivial-topological-vector-bundle]{trivial topological \(k\)-vector bundle} \(X \times k^{n}\) corresponds to \((\mathcal{C}_{X}(k))^{\oplus n}\).
    \item A \hyperref[def:section-topological-vector-bundle]{section} of \(E \xrightarrow{\pi} X\) corresponds to a \hyperref[def:global-section-sheaf]{global section} of \(\mathcal{C}_{X}(k)\).
  \end{itemize}
\end{theorem}
\begin{proof}
  In \autoref{prop:sheaf-of-sections-topological-vector-bundle-is-C-x-module} we have seen that \(\mathcal{E}\) defines a \hyperref[def:functor]{functor}
  \begin{equation*}
    \Vect_{n}(X) \to \category{LocFree}_{n}(\mathcal{C}_{X}(k))
  \end{equation*}
  Therefore we only need to construct the inverse correspondence.

  Let \(\mathcal{E}\) be a \hyperref[def:locally-free-sheaf]{locally free sheaf of \(\mathcal{C}_{X}(k)\)-modules}.
  Then there exists an open covering \(\mathcal{U} = \{U_{i}\}_{i \in \mathcal{I}}\) of \(X\) such that for every \(U_{i} \in \mathcal{U}\) there is an \hyperref[def:categorical-isomorphism]{isomorphism} of \hyperref[def:sheaf-of-modules]{sheaves of \(\mathcal{C}_{X}(k)\)-modules}
  \begin{equation*}
    \tau|_{i}: \mathcal{E}|_{U_{i}} \isoarrow \left( \mathcal{C}_{X}(k)|_{U_{i}} \right)^{\oplus n}
  \end{equation*}
  Therefore on \(U_{i} \cap U_{j}\) we have an isomorphism of \hyperref[def:locally-free-sheaf]{locally free sheaves of \(\mathcal{C}_{X}(k)\)-modules}
  \begin{equation*}
    \tau_{j} \circ \tau_{i}^{-1}: \left( \mathcal{C}_{X}(k)|_{U_{i} \cap U_{j}}  \right)^{\oplus n} \isoarrow \left( \mathcal{C}_{X}(k)|_{U_{i} \cap U_{j}}  \right)^{\oplus n}
  \end{equation*}
  Note that this isomorphism is given by an \(n \times n\) matrix \(M_{i,j}\) consisting of elements of \(\mathcal{C}_{X}(k)|_{U_{i} \cap U_{j}}\) that is invertible at every point.

  This matrices are not arbitrary.
  Indeed they satisfy the \textit{cocycle condition}:
  \begin{equation*}
    M_{i, j} \cdot M_{j, k} \cdot M_{k, i} = \operatorname{Id}_{n}
  \end{equation*}
  
  Therefore we can assign to \(\mathcal{E}\) the \hyperref[def:vector-bundle]{vector bundle} of rank \(n\), \hyperref[ex:trivial-vector-bundle]{trivial} on the open sets \(U_{i}\) of the covering, and whose transitions are the matrics \(M_{i,j}\).
  Less abstractly, \(E\) is obtained by gluing \(U_{i} \times \mathbb{R}^{n}\) via the identification
  \begin{equation*}
    U_{i} \cap U_{j} \times \mathbb{R}^{n} \subset U_{i} \times \mathbb{R}^{n} \isoarrow U_{i} \times U_{j} \times \mathbb{R}^{n} \subset U_{j} \times \mathbb{R}^{n}
  \end{equation*}
  given by \(\id \times M_{i, j}\).

  This defines a \hyperref[def:functor]{functor}
  \begin{equation*}
    \category{LocFree}_{n}(\mathcal{C}_{X}(k)) \to \Vect_{n}(X)
  \end{equation*}
  that by construction is \textit{inverse} to the previous \hyperref[def:functor]{functor}.
\end{proof}


\begin{definition}[Smooth vector bundle]\label{def:smooth-vector-bundle}
  Let \(M\) be a smooth manifold.
  Let \(k \in \{\mathbb{R}, \mathbb{C}\}\) be either the field of real or complex numbers.

  A \textbf{smooth \(k\)-vector bundle} is a \hyperref[def:topological-vector-bundle]{topological \(k\)-vector bundle} \(E \xrightarrow{\pi} M\) such that
  \begin{enumerate}
    \item \(E\) is a smooth manifold.
    \item The projection map \(\pi: E \to M\) is smooth.
    \item The local trivialization maps \(\phi_{i}: U_{i} \times k^{n_{i}} \isoarrow \pi^{-1}(U_{i})\) are diffeomorphism.
  \end{enumerate}

  
  A \textbf{morphism} of smooth \(k\)-vector bundles from \(E_{1} \xrightarrow{\pi_{1}} M\) to \(E_{2} \xrightarrow{\pi_{2}} M\) is \hyperref[def:topological-vector-bundle]{morphism of topological vector bundles} \(f: E_{1} \to E_{2}\) such that \(f\) is smooth.

  We denote by \(\Vect_{k}^{\infty}(X)\) the category of smooth \(k\)-vector bundles over \(X\).
\end{definition}

\begin{definition}[Holomorphic vector bundle]\label{def:holomorphic-vector-bundle}
  Let \(M\) be a complex manifold.

  A \textbf{holomorphic vector bundle} is a \hyperref[def:topological-vector-bundle]{topological \(\mathbb{C}\)-vector bundle} \(E \xrightarrow{\pi} M\) such that
  \begin{enumerate}
    \item \(E\) is a complex manifold.
    \item The projection map \(\pi: E \to M\) is an holomorphism.
    \item The local trivialization maps \(\phi_{i}: U_{i} \times \mathbb{C}^{n_{i}} \isoarrow \pi^{-1}(U_{i})\) are biholomorphisms\footnote{A map \(f: \mathbb{C} \to \mathbb{C}\) is biholomorphic is it is holomorphic, bijective and its inverse is holomorphic.}.
  \end{enumerate}

  
  A \textbf{morphism} of holomorphic vector bundles from \(E_{1} \xrightarrow{\pi_{1}} M\) to \(E_{2} \xrightarrow{\pi_{2}} M\) is \hyperref[def:topological-vector-bundle]{morphism of topological vector bundles} \(f: E_{1} \to E_{2}\) such that \(f\) is holomorphic.

  We denote by \(\Vect_{\mathbb{C}}^{\omega}(X)\) the category of holomorphic vector bundles over \(X\).
\end{definition}

\begin{remark}
  \autoref{thm:topological-vector-bundle-is-free-sheaf} can be proved analogously to the case of \hyperref[def:smooth-vector-bundle]{smooth} and \hyperref[def:holomorphic-vector-bundle]{holomorphic} vector bundles.
\end{remark}


\end{document}
